\thesisTitle{EMPIRICAL PERFORMANCE ANALYSIS OF OPERATOR FUSION ON NEURAL PROCESSING UNITS (NPUs): A STUDY ON TIME AND I/O COMPLEXITY}
\thesisTitleTR{YAPAY ZEKA İŞLEM BİRİMLERİNDE (NPU) OPERATÖR FÜZYONUNUN AMPİRİK PERFORMANS ANALİZİ: ZAMAN VE I/O KARMAŞIKLIĞI ÜZERİNE BİR İNCELEME}
%\thesisYear{2026}
%\thesisMonth{April}\thesisMonthTR{Nisan}
\thesisDefenseDay{12}  % day of defense, 1, 2, ..., 31
\thesisDegree{MSc}  %% PhD || MSc
\thesisDegreeTR{Yüksek Lisans} %% Doktora || Yüksek Lisans
\thesisProgram{Software Engineering}
\thesisProgramTR{Yazılım Mühendisliği}
\thesisAuthor{Yusuf Talha ARABACI}
\thesisAdvisorTitle{Dr. Öğr. Üyesi}
\thesisAdvisorTitleTR{Dr. Öğr. Üyesi}
\thesisAdvisor{Emrullah DEMİRAL}
\thesisAdvisorDepartment{Department of Software Engineering}

% Template customization for course project (overrides defaults in the class)
\thesisDocTypeEN{COURSE PROJECT REPORT}
\thesisDocTypeTR{DERS PROJESİ RAPORU}
\thesisAdvisorLabelEN{Project Advisor}
\thesisAdvisorLabelTR{Proje Danışmanı}
\thesisPreparedAsEN{Course Project Report}
\thesisPreparedAsTR{Ders Projesi Raporu}
\thesisCourseNameEN{Advanced Algorithm Analysis -- Course Project}
\thesisCourseNameTR{İleri Algoritma Analizi -- Ders Projesi}

\thesisChairman{(Jüri Başkanı)}
\thesisMemberTWO{(Üye 1)}
\thesisMemberTHREE{(Üye 2)}
\thesisMemberFOUR{}
\thesisMemberFIVE{}

\thesisPages{29}

\thesisAbstract{ 
Traditional algorithm analysis generally ignores hardware memory constraints (RAM Model). However, in modern NPU architectures, the primary performance loss stems from data transfer to and from memory rather than computation. This study examines the impact of graph optimization levels provided by the ONNX Runtime compiler on execution time. Empirical tests reveal that reducing the number of operations by merging them (Fusion), even without changing the algorithm's asymptotic growth class, significantly accelerates inference times by lowering constant operation costs.
}

\thesisKeywords{ONNX Runtime, graph optimization, operator fusion, NPU, OpenVINO, memory wall, I/O complexity, performance benchmarking.}

\thesisOzet{ 
Geleneksel algoritma analizi genellikle donanım bellek kısıtlarını göz ardı eder (RAM Modeli). Ancak günümüz NPU mimarilerinde asıl performans kaybı hesaplamadan değil, verinin belleğe gidip gelmesinden kaynaklanmaktadır. Bu çalışma, ONNX Runtime derleyicisinin sunduğu graf optimizasyon seviyelerinin çalışma zamanına etkisini incelemektedir. Yapılan ampirik testler, işlemleri birleştirerek sayısını azaltmanın (Fusion), algoritmanın asimptotik büyüme sınıfını değiştirmese bile, sabit işlem maliyetlerini düşürerek çıkarım (inference) sürelerini net bir şekilde hızlandırdığını ortaya koymaktadır.
}

\thesisKeywordsTR{ONNX Runtime, grafik optimizasyonu, operatör füzyonu, NPU, OpenVINO, profiling, performans kıyaslama.}

\thesisScienceCode{0}

\thesisAcknowledgment{Bu çalışmanın yürütülmesinde desteği olan ders yürütücüsüne ve katkı sağlayan herkese teşekkür ederim.}

\thesisResume{Yusuf Talha ARABACI (yusuftalhaarabaci@hotmail.com) -- Yazılım Mühendisliği Yüksek Lisans öğrencisi. İlgi alanları: derleyici optimizasyonları, grafik optimizasyonları, hızlandırıcılar ve performans analizi.}

\thesisSymbolList{
\thesisSymbol{$S$}{Speedup (baseline/optimized)}
\thesisSymbol{$t$}{Inference latency (ms)}
\thesisSymbol{$N$}{Graph node count}
\thesisSymbol{$AI$}{Arithmetic intensity (FLOP/byte)}
}

\thesisAbbrList{
\thesisAbbr{ORT}{ONNX Runtime}
\thesisAbbr{EP}{Execution Provider}
\thesisAbbr{NPU}{Neural Processing Unit}
\thesisAbbr{CPU}{Central Processing Unit}
\thesisAbbr{OV}{OpenVINO}
\thesisAbbr{FLOP}{Floating-point operation}
}

\makeThesisPreamble
\pagenumbering{arabic}